\examplegloss{residual-action}{Residual action}{
\begin{enumerate}[label=(\arabic*),leftmargin=*,itemsep=2pt,topsep=3pt]
\item On $\PP(2,1)$ the action $r\cdot[x:y]=[rx:y]$ commutes with the cutting action. On the interior coordinate $z=x/y^2$ it becomes $z\mapsto rz$. The point $z=0$ is fixed by this residual action, but has trivial stack stabiliser.
\item The boundary $[1:0]$ is also fixed as an isomorphism class under the residual action. Its stack stabiliser is $\mu_2$, coming from $t^2=1$ in the cutting group. These automorphisms and the residual symmetry are different data.
\end{enumerate}}

\examplegloss{residual-gerbe-at-a-closed-point}{Residual gerbe at a closed point}{
\begin{enumerate}[label=(\arabic*),leftmargin=*,itemsep=2pt,topsep=3pt]
\item At the origin of $[\A^1_z/\mu_2]$, the reduced closed substack is $B\mu_2$: it retains one point together with its identity and sign automorphisms. This is the residual gerbe of the origin.
\item At a nonzero coarse point of the same quotient, the scalar action is free. The residual gerbe is an ordinary $\Spec k$. The two cases differ through their automorphisms, although both coarse points are ordinary points of $\A^1$.
\end{enumerate}}

\examplegloss{root-divisor}{Root divisor}{
\begin{enumerate}[label=(\arabic*),leftmargin=*,itemsep=2pt,topsep=3pt]
\item On $\sqrt[2]{(\A^1_u,\{0\})}=[\A^1_z/\mu_2]$, the root divisor is cut out by $z=0$, so it is $B\mu_2$. Since $u=z^2$, the pullback of the ordinary origin divisor is twice this divisor.
\item On $\PP(2,1)$, the section $y$ vanishes simply along the root divisor $\cD=\{y=0\}$. The coarse boundary point has local parameter $y^2/x$, which vanishes to order two upstairs. Thus again $\pi^*D=2\cD$.
\end{enumerate}}

\examplegloss{root-gerbe}{Root gerbe}{
\begin{enumerate}[label=(\arabic*),leftmargin=*,itemsep=2pt,topsep=3pt]
\item Square roots of the trivial line on $\Spec k$ form $B\mu_2$. The root is a line with a specified square trivialisation. Multiplication by $-1$ preserves that trivialisation, giving a nonidentity automorphism.
\item Square roots of $\OO_{\A^1}$ form $\A^1\times B\mu_2$. Every geometric point, including nonzero points, has a $\mu_2$ stabiliser. This illustrates why omitting a root section leaves stabilisers everywhere, rather than only over a divisor.
\end{enumerate}}

\examplegloss{root-stack-of-a-divisor}{Root stack of a divisor}{
\begin{enumerate}[label=(\arabic*),leftmargin=*,itemsep=2pt,topsep=3pt]
\item Taking the square root of the origin divisor in $\A^1_u$ gives $[\A^1_z/\mu_2]$ with $u=z^2$. Away from zero the root section is nonzero and removes all automorphisms; at zero its sign ambiguity remains.
\item Taking the square root of one point $D\subset\PP^1$ gives $\PP(2,1)$. Its ordinary affine chart away from $D$ and cyclic quotient chart near $D$ show explicitly how a root stack modifies just the chosen divisor.
\end{enumerate}}

\examplegloss{scheme-theoretic-fibre}{Scheme-theoretic fibre}{
\begin{enumerate}[label=(\arabic*),leftmargin=*,itemsep=2pt,topsep=3pt]
\item For the scheme map $\A^1_z\to\A^1_u$, $u=z^2$, the fibre over zero is $\Spec k[z]/(z^2)$. It has one underlying point but retains the nonzero nilpotent class of $z$.
\item For the stack coarse map $[\A^1_z/\mu_2]\to\A^1_u$, base change gives $[\Spec k[z]/(z^2)/\mu_2]$. Its reduction is $B\mu_2$. Merely listing its geometric objects would miss the thickening.
\end{enumerate}}

\examplegloss{section-ring}{Section ring}{
\begin{enumerate}[label=(\arabic*),leftmargin=*,itemsep=2pt,topsep=3pt]
\item For $M=\PP^1$ and $L=\OO(1)$, sections of $L^m$ are degree-$m$ homogeneous polynomials in $X_0,X_1$. Thus the section ring is $k[X_0,X_1]$, with both variables in degree one.
\item For $M=\A^1_x$ and $L=\OO_M$, introduce its constant degree-one section $h$. Then $R=k[x,h]$, with $\deg x=0$ and $\deg h=1$. The nonconstant function $x$ belongs to the degree-zero algebra.
\end{enumerate}}

\examplegloss{section-weight}{Section weight}{
\begin{enumerate}[label=(\arabic*),leftmargin=*,itemsep=2pt,topsep=3pt]
\item If $t\cdot x=t^2x$ on $\A^1$ and $L=\OO$ has trivial fibre action, direct pullback gives $x(tp)=t^2x(p)$. The section $x$ has weight $2$, and $x^3$ has weight $6$.
\item Over a fixed point, give a line geometric fibre action $v\mapsto tv$. The section action includes the inverse fibre map, so its constant nonzero section has weight $-1$. This is the sign used for the level-one twist.
\end{enumerate}}

\examplegloss{semiprojective}{Semiprojective}{
\begin{enumerate}[label=(\arabic*),leftmargin=*,itemsep=2pt,topsep=3pt]
\item For $M=\A^1$, the affinisation is the identity $\A^1\to\Spec k[x]$. An isomorphism is projective, so $\A^1$ is semiprojective in the convention here, although it is not proper over $k$.
\item For $M=\PP^1$, all global regular functions are constant. Its affinisation is therefore the projective morphism $\PP^1\to\Spec k$. These two examples show how the convention includes both affine and projective varieties.
\end{enumerate}}

\examplegloss{semistable-point}{Semistable point}{
\begin{enumerate}[label=(\arabic*),leftmargin=*,itemsep=2pt,topsep=3pt]
\item In the lower cut with point weights $(2,1)$ and level $1$, $(0,1)$ is semistable because $y\one_1$ is nonzero there. The point $(1,0)$ is semistable because $x\one_1^2$ is nonzero there.
\item At level $0$, the constant twist section has positive tensor degree and weight zero, so it is nonzero even at $(0,0)$. That point is semistable although its stabiliser is the full group $\Gm$.
\end{enumerate}}

\examplegloss{separated-stack}{Separated stack}{
\begin{enumerate}[label=(\arabic*),leftmargin=*,itemsep=2pt,topsep=3pt]
\item The scheme $\A^1$, viewed as a stack, is separated: its diagonal in $\A^1\times\A^1$ is the closed subscheme $x=y$. A closed immersion is proper.
\item The stack $B\mu_2$ is also separated. After trivialising two torsors, its isomorphism sheaf is a copy of the finite group $\mu_2$. Its diagonal is finite and therefore proper, despite its nontrivial automorphisms.
\end{enumerate}}

\examplegloss{smooth-stack}{Smooth stack}{
\begin{enumerate}[label=(\arabic*),leftmargin=*,itemsep=2pt,topsep=3pt]
\item The stack $[\A^1/\mu_2]$ is smooth because $\A^1$ is a smooth étale atlas. Its coarse space is also a smooth affine line, with coordinate $z^2$.
\item The stack $[\A^2/\mu_2]$ for simultaneous sign change is smooth for the same atlas reason. Its coarse space is the singular surface $ac=b^2$, with $a=x^2$, $b=xy$, $c=y^2$. Smoothness of the stack does not force smoothness of its coarse space.
\end{enumerate}}

\examplegloss{stabiliser-or-isotropy-group}{Stabiliser or isotropy group}{
\begin{enumerate}[label=(\arabic*),leftmargin=*,itemsep=2pt,topsep=3pt]
\item At $z=0$ in $[\A^1_z/\mu_2]$, both $1$ and $-1$ fix the representative. The stabiliser is $\mu_2$, and it acts by sign on the tangent direction.
\item At any point of $\PP(2,1)$ with $y\ne0$, a stabilising $t$ must satisfy $ty=y$, hence $t=1$. Its stabiliser is trivial even though the group used in the quotient presentation is one-dimensional.
\end{enumerate}}

\examplegloss{stable-point}{Stable point}{
\begin{enumerate}[label=(\arabic*),leftmargin=*,itemsep=2pt,topsep=3pt]
\item For scalar multiplication on $\A^2\setminus0$, each orbit is a line with the origin removed. It is closed within this punctured plane and has trivial stabiliser. These points give the usual stable presentation of $\PP^1$.
\item For weights $(2,1)$, the orbit of $(1,0)$ is the nonzero $x$-axis. It is closed in $\A^2\setminus0$ and has stabiliser $\mu_2$. Thus this point is stable, although the action there is not free.
\end{enumerate}}

\examplegloss{stack}{Stack}{
\begin{enumerate}[label=(\arabic*),leftmargin=*,itemsep=2pt,topsep=3pt]
\item The scheme $\A^1$ gives a stack whose objects over $T$ are regular functions on $T$ and whose only arrows are identities. Compatible local functions glue, so ordinary schemes fit the stack axioms.
\item The stack $B\mu_2$ instead assigns to $T$ its groupoid of $\mu_2$-torsors. These torsors and their isomorphisms glue, while geometric objects retain sign automorphisms. The same gluing principle accommodates nontrivial symmetries.
\end{enumerate}}

\examplegloss{tame-stack}{Tame stack}{
\begin{enumerate}[label=(\arabic*),leftmargin=*,itemsep=2pt,topsep=3pt]
\item On $B\mu_2$, a vector space with involution $\sigma$ is pushed to its invariant part. The averaging operator $v\mapsto(v+\sigma v)/2$ makes taking invariants exact. This exhibits tameness over the coarse point.
\item On $[\A^1/\mu_2]$, the same averaging works for equivariant modules over $k[z]$. Their invariant modules form the coarse pushforward over $k[z^2]$, and exact sequences remain exact. Tameness does not remove the sign characters on stack line fibres.
\end{enumerate}}

\examplegloss{tensor-degree}{Tensor degree}{
\begin{enumerate}[label=(\arabic*),leftmargin=*,itemsep=2pt,topsep=3pt]
\item On $\PP^1$, the monomial $X_0^2$ is a section of $\OO(1)^{\otimes2}=\OO(2)$, so its tensor degree is $2$. This degree counts powers of the line bundle.
\item On the affine-line example with $t\cdot x=t^2x$, the element $x^7h^3$ lies in $H^0(\A^1,\OO^{\otimes3})$. Its tensor degree is $3$, while its group weight is $14$. These are independent indices.
\end{enumerate}}

\examplegloss{topology-in-descent-theory}{Topology, in descent theory}{
\begin{enumerate}[label=(\arabic*),leftmargin=*,itemsep=2pt,topsep=3pt]
\item The two standard affine charts of $\PP^1$ form a Zariski cover. A line bundle can be described by trivial lines on these charts and a transition function on their overlap.
\item The squaring map $\Gm\to\Gm$ is an étale cover. It allows a square root of the target coordinate after base change, although no such root exists on a nonempty Zariski open. Changing the topology changes which local choices are available for gluing.
\end{enumerate}}

\examplegloss{torsor}{Torsor}{
\begin{enumerate}[label=(\arabic*),leftmargin=*,itemsep=2pt,topsep=3pt]
\item The product $\Gm\times T\to T$, with multiplication on the first factor, is a trivial $\Gm$-torsor. Its identity section gives the chosen trivialisation; multiplying that section by a unit gives another one.
\item The map $\Spec k[s,s^{-1}]\to\Spec k[t,t^{-1}]$, $t=s^2$, is a $\mu_2$-torsor. It has no global section because $t$ has no square root in $k[t,t^{-1}]$, but it becomes trivial after pulling back along itself.
\end{enumerate}}

\examplegloss{torus}{Torus}{
\begin{enumerate}[label=(\arabic*),leftmargin=*,itemsep=2pt,topsep=3pt]
\item The group $\Gm$ is a one-dimensional torus. It acts on $\A^1$ by scalar multiplication, or by the weighted formula $x\mapsto t^2x$. Its characters are indexed by integers.
\item The product $\Gm\times\Gm$ is a two-dimensional torus, acting on $\A^2$ by $(t_1,t_2)\cdot(x,y)=(t_1x,t_2y)$. A character has the form $t_1^a t_2^b$, indexed by the pair $(a,b)\in\mathbb Z^2$.
\end{enumerate}}

\examplegloss{transporter}{Transporter}{
\begin{enumerate}[label=(\arabic*),leftmargin=*,itemsep=2pt,topsep=3pt]
\item For the sign action on $\A^1$, the transporter from $1$ to $-1$ is $\{-1\}$. It is the space of arrows between these two representatives of the same quotient object.
\item The transporter from $0$ to itself is all of $\mu_2$, because both group elements fix zero. Thus the same transporter construction gives isomorphisms between different representatives and automorphisms when the representatives coincide.
\end{enumerate}}

\examplegloss{unramified-morphism}{Unramified morphism}{
\begin{enumerate}[label=(\arabic*),leftmargin=*,itemsep=2pt,topsep=3pt]
\item The map $\mathbb G_{m,s}\to\mathbb G_{m,t}$, $t=s^2$, is unramified: the relative differential relation $2s\,ds=0$ forces $ds=0$ because $2s$ is invertible. This map is also flat, hence étale.
\item The closed inclusion $\{0\}\hookrightarrow\A^1$ is unramified, since a point has no relative differentials over $\A^1$. It is not flat and is not étale. Thus unramified alone does not mean that a morphism is an étale chart.
\end{enumerate}}

\examplegloss{vector-bundle}{Vector bundle}{
\begin{enumerate}[label=(\arabic*),leftmargin=*,itemsep=2pt,topsep=3pt]
\item The sheaf $\OO_{\A^1}^{\oplus2}$ is the trivial rank-two vector bundle. Its fibre at every point is $k^2$, and its two standard basis sections give a global frame.
\item On $B\mu_2$, take the representation $\sigma(v_1,v_2)=(v_1,-v_2)$. It gives a rank-two vector bundle with one trivial and one sign summand. Its nontrivial fibre action prevents descent to a vector bundle on the coarse point.
\end{enumerate}}

\examplegloss{veronese-subring}{Veronese subring}{
\begin{enumerate}[label=(\arabic*),leftmargin=*,itemsep=2pt,topsep=3pt]
\item For $k[u,v]$ with both variables in degree one, the second Veronese is generated by $u^2,uv,v^2$. After regrading these have degree one and satisfy $(uv)^2=u^2v^2$. Its projective spectrum is still $\PP^1$.
\item For $k[A,B]$ with degrees $(2,1)$, every even-degree monomial is a product of $A$ and $B^2$. The second Veronese is therefore $k[A,B^2]$, with both generators in regraded degree one, giving the coarse projective line of the toy cut.
\end{enumerate}}

\examplegloss{weight}{Weight}{
\begin{enumerate}[label=(\arabic*),leftmargin=*,itemsep=2pt,topsep=3pt]
\item The point action $x\mapsto t^2x$ gives tangent weight $2$ at the origin. With the direct-pullback action used here, the coordinate function $x$ also has weight $2$.
\item Give a line over a fixed point the geometric fibre action $v\mapsto t^3v$. Its fibre weight is $3$, but its constant section has weight $-3$ under the declared section convention. Naming the representation is essential when stating a weight.
\end{enumerate}}

\examplegloss{weighted-projective-stack}{Weighted projective stack}{
\begin{enumerate}[label=(\arabic*),leftmargin=*,itemsep=2pt,topsep=3pt]
\item With weights $(1,1)$ on $\A^2\setminus0$, the action is free and the quotient stack $\PP(1,1)$ is the ordinary projective line. Either nonzero coordinate can be normalised without a finite ambiguity.
\item With weights $(2,1)$, the quotient is $\PP(2,1)$. The chart $y\ne0$ is $\A^1$, while the boundary has stabiliser $\mu_2$. Its coarse space is again $\PP^1$, so the weights change stack geometry without changing this coarse curve.
\end{enumerate}}

\examplegloss{weight-space}{Weight space}{
\begin{enumerate}[label=(\arabic*),leftmargin=*,itemsep=2pt,topsep=3pt]
\item For $k[x]$ with $x\mapsto t^2x$ under direct pullback, the weight-$2a$ space is $k x^a$ for $a\ge0$. Every odd-weight space is zero. The invariant part is the weight-zero space $k$.
\item For the $\mu_2$-representation $k^2$ with action $(v_1,v_2)\mapsto(v_1,-v_2)$, the two coordinate axes are the trivial and sign weight spaces. Here the weight labels are residues modulo $2$.
\end{enumerate}}

\examplegloss{zariski-versus-etale-local}{Zariski versus étale local}{
\begin{enumerate}[label=(\arabic*),leftmargin=*,itemsep=2pt,topsep=3pt]
\item On $\mathbb G_{m,t}$, an equation $s^2=t$ has no solution on any nonempty Zariski open: it has no solution even in $k(t)$. Pulling back by the étale squaring cover supplies exactly such a root.
\item On the chart $x\ne0$ of $\PP(2,1)$, setting $x=1$ requires choosing $t$ with $t^2=x^{-1}$, which is possible étale locally and leaves a sign ambiguity. On $y\ne0$, setting $y=1$ instead uses the unique regular choice $t=y^{-1}$.
\end{enumerate}}
